\documentclass[../../main.tex]{subfiles}
\begin{document}
\section{Lie 对应}\label{section:ch20-lie-correspondence}

我们关于 Lie 群的许多结果表明，Lie 群的本质性质如何反映到其 Lie 代数中，反之亦然。这自然引出一个问题：Lie 群与 Lie 代数之间（至少在其同构类之间）的对应在多大程度上是一一对应的？

我们在第 8 章已经看到，把 Lie 群映到其 Lie 代数、把 Lie 群同态映到其诱导 Lie 代数同态的指派是从 Lie 群范畴到有限维 Lie 代数范畴的一个函子。由于函子把同构映到同构，可知同构的 Lie 群有同构的 Lie 代数。反过来易见不成立：$\mathbb{R}^n$ 与 $\mathbb{T}^n$ 都有 $n$ 维交换 Lie 代数，它们显然彼此同构，但 $\mathbb{R}^n$ 与 $\mathbb{T}^n$ 绝非同构的 Lie 群。然而，正如我们将在本节看到的，若把注意力限于单连通 Lie 群，我们确实得到一一对应。

为证明这一对应，我们需要在给定两个单连通群的代数之间的同构时，构造二者之间同构的方法。定理\cref{theorem:ch08-44} 表明每个 Lie 群同态都给出一个 Lie 代数同态。利用 Lie 代数作用基本定理，我们可证明下列部分逆命题。

\begin{theorem}\label{theorem:ch20-同态提升}
设 $G$ 与 $H$ 是 Lie 群，$G$ 单连通，$\mathfrak g$ 与 $\mathfrak h$ 是它们的 Lie 代数。对任意 Lie 代数同态 $\varphi\colon\mathfrak g\to\mathfrak h$，存在唯一的 Lie 群同态 $\Phi\colon G\to H$，使 $\Phi_*=\varphi$。
\end{theorem}

\begin{proof}
Lie 代数同态 $\varphi\colon\mathfrak g\to\mathfrak h\subseteq\mathfrak X(H)$ 特别地是 $H$ 上的一个完全 $\mathfrak g$-作用（因为每个左不变向量场都是完全的）。于是由定理\cref{theorem:ch20-Lie代数作用基本定理}，存在唯一的光滑右 $G$-作用 $\sigma\colon H\times G\to H$，使 $\varphi$ 为其无穷小生成元。用记号 $\widetilde X=\varphi(X)$（$X\in\mathfrak g$），并用 $h\cdot g=\sigma(h,g)$（$h\in H$，$g\in G$）。定义光滑映射 $\Phi\colon G\to H$，$\Phi(g)=e\cdot g$（$e$ 为 $H$ 的单位元）。我们将证 $\Phi$ 即所求同态。

引理\cref{lemma:ch20-轨道映射相关} 表明对每个 $h\in H$、每个 $X\in\mathfrak g$，向量场 $X$ 与 $\widetilde X$ 是 $\sigma_h$-相关的。由命题\cref{proposition:ch09-06}，$\sigma_h$ 把 $X$ 的积分曲线映为 $\widetilde X$ 的积分曲线。因此 $t\mapsto h\cdot\exp tX$ 是从 $h$ 出发的 $\widetilde X$ 的积分曲线。另一方面，由于 $\widetilde X$ 是 $H$ 上的左不变向量场，$H$ 中的左平移把 $\widetilde X$ 的积分曲线映为 $\widetilde X$ 的积分曲线。对任意 $h\in H$、$X\in\mathfrak g$，取 $h=e\cdot g$（$g\in G$）应用此结论，得
\begin{align*}
(e\cdot g)\cdot(e\cdot\exp tX)=(e\cdot g)\exp tX=e\cdot(g\exp tX).
\end{align*}
（最后一等式由 $\sigma$ 是作用得到。）用 $\Phi$ 改写，即
\begin{align*}
\Phi(g)\Phi(\exp tX)=\Phi(g\exp tX).
\end{align*}
由于 $G$ 连通，由命题\cref{proposition:ch07-14} 它由指数映射的像生成，故这蕴含 $\Phi$ 是同态。

为看出 $\Phi_*=\varphi$，任取 $X\in\mathfrak g$。$\varphi$ 是 $\sigma$ 的无穷小生成元意味着
\begin{align*}
\varphi(X)_e=\frac{\mathrm{d}}{\mathrm{d}t}\bigg|_{t=0}(e\cdot\exp tX)=\frac{\mathrm{d}}{\mathrm{d}t}\bigg|_{t=0}\Phi(\exp tX)=\mathrm{d}\Phi_e(X_e).
\end{align*}
由于 $\Phi$ 由 $\mathrm{d}\Phi_e$ 的作用决定，这蕴含 $\Phi_*X=\varphi(X)$。

证明通过援引习题 20-17 完成，它表明 $\Phi$ 是具有此性质的唯一同态。
\end{proof}

\begin{corollary}\label{corollary:ch20-单连通同构}
若 $G$ 与 $H$ 是单连通 Lie 群且其 Lie 代数同构，则 $G$ 与 $H$ 同构。
\end{corollary}

\begin{proof}
设 $\mathfrak g,\mathfrak h$ 分别为 $G,H$ 的 Lie 代数，$\varphi\colon\mathfrak g\to\mathfrak h$ 为二者间的 Lie 代数同构。由上述定理，存在 Lie 群同态 $\Phi\colon G\to H$ 与 $\Psi\colon H\to G$ 满足 $\Phi_*=\varphi$、$\Psi_*=\varphi^{-1}$。恒等映射 $\operatorname{Id}_G$ 与复合 $\Psi\circ\Phi$ 都是从 $G$ 到自身的 Lie 群同态，其诱导 Lie 代数同态都等于恒等，故定理\cref{theorem:ch20-同态提升} 的唯一性部分蕴含 $\Psi\circ\Phi=\operatorname{Id}_G$。同理 $\Phi\circ\Psi=\operatorname{Id}_H$，故 $\Phi$ 是 Lie 群同构。
\end{proof}

现在我们准备好证明主定理。

\begin{theorem}{[Lie 对应]}\label{theorem:ch20-Lie对应}
有限维 Lie 代数的同构类与单连通 Lie 群的同构类之间存在一一对应，其方式是把每个单连通 Lie 群与其 Lie 代数相伴。
\end{theorem}

\begin{proof}
需证把单连通 Lie 群映到其 Lie 代数的函子在同构意义下既满又单。单射恰为推论\cref{corollary:ch20-单连通同构} 的内容。

为证满射，设 $\mathfrak g$ 为任意有限维 Lie 代数。由 Lie 的基本定理（定理\cref{theorem:ch20-Lie基本定理}），存在一个单连通 Lie 群，其 Lie 代数同构于 $\mathfrak g$。
\end{proof}

\begin{remark}{[Lie 的基本定理]}\label{remark:ch20-Lie基本定理历史}
正如上述定理名称所提示，Lie 对应定理的一个版本是在十九世纪由 Sophus Lie 证明的。然而，由于流形与单连通性等整体拓扑概念当时尚未提出，他所能证明的本质上是该定理的局部版本。

Lie 不是把 Lie 群作为抽象对象来研究，而是研究 Euclid 空间开子集上的向量场及其生成的（局部）群作用。定义\emph{局部 Lie 群}为某个有限维向量空间 $V$ 中的开子集 $U$，连同元素 $e\in U$ 与光滑映射 $m\colon U\times U\to V$（乘法）及 $i\colon U\to V$（取逆），满足对所有充分接近 $e$、使两边都有定义的 $g,h,k$：
\begin{align*}
m\bigl(g,m(h,k)\bigr)=m\bigl(m(g,h),k\bigr),\quad m(e,g)=g=m(g,e),\quad m\bigl(i(g),g\bigr)=e=m\bigl(g,i(g)\bigr).
\end{align*}
左平移 $L_g\colon U\to V$ 与普通 Lie 群同样定义，向量场 $X\in\mathfrak X(U)$ 称为左不变的，若对所有使 $m(g,g')\in U$ 的 $g,g'$ 有 $\mathrm{d}(L_g)_{g'}(X_{g'})=X_{m(g,g')}$。两个局部 Lie 群称为局部同构，若存在从 $U$ 中 $e$ 的邻域到 $U'$ 中 $e'$ 的邻域的微分同胚，在相应运算有定义时把 $e$ 映到 $e'$、$m$ 映到 $m'$、$i$ 映到 $i'$。局部 Lie 群在开子集 $W\subseteq\mathbb{R}^n$ 上的局部（左或右）作用与普通作用同样定义，只是 $gx$（或 $xg$）只需在 $U\times W$ 中 $\{e\}\times W$ 的一个邻域上有定义。任意 Lie 群中单位元的一个坐标邻域是局部 Lie 群，而 Lie 群在光滑流形上的任意光滑作用限制到任意充分小的坐标邻域上即为局部作用。
\end{remark}

\begin{theorem}{[Sophus Lie 基本定理]}\label{theorem:ch20-Lie基本定理}
每个有限维 Lie 代数都是某个局部 Lie 群的 Lie 代数；两个局部 Lie 群局部同构当且仅当其 Lie 代数同构。
\end{theorem}
\end{document}
